\chapter{Reports}
\label{ch:reports}

\section{Why reports at all}

A dashboard answers questions while the user is looking at it. A report is
something that can be filed, sent, or read a year later. The application
produces both formats people actually use: Excel for anyone who wants to poke at
the numbers, and PDF for anything that gets read rather than edited.

\section{Excel}

The Excel export is built with \code{openpyxl}. All the styling lives in one
function, so every sheet looks the same:

\begin{lstlisting}[style=py]
def _write_sheet(ws, headers: list[str], rows: list[list]) -> None:
    for c, h in enumerate(headers, 1):
        cell = ws.cell(row=1, column=c, value=h)
        cell.fill = _HEADER_FILL
        cell.font = _HEADER_FONT
    for r, row in enumerate(rows, 2):
        for c, val in enumerate(row, 1):
            cell = ws.cell(row=r, column=c, value=val)
            if r % 2 == 0:
                cell.fill = _ALT_FILL
            if isinstance(val, (int, float, Decimal)) and c > 1:
                cell.number_format = _money_format()
    for col in ws.columns:
        width = max((len(str(c.value)) for c in col if c.value is not None), default=8)
        ws.column_dimensions[get_column_letter(col[0].column)].width = width + 4
\end{lstlisting}

Two details matter. Amounts are written as numbers with a currency number
format, not as pre-formatted strings, so the spreadsheet can still sum them. And
the column width is computed from the longest value, which removes the usual
first action of every reader, widening the columns by hand.

\section{PDF}

The PDF reports use \code{reportlab} for the document and \code{matplotlib} for
the charts, which are rendered to PNG in memory and embedded.

\begin{lstlisting}[style=py]
def _save(fig) -> bytes:
    buf = io.BytesIO()
    fig.savefig(buf, format="png", dpi=130, bbox_inches="tight", facecolor="white")
    plt.close(fig)
    buf.seek(0)
    return buf.getvalue()
\end{lstlisting}

Nothing touches the filesystem. The chart exists as bytes, goes into the
document, and disappears. On a hosting platform with a read-only or temporary
filesystem that is the only approach that works, and it avoids leaving files
behind after a failed request.

The \code{Agg} backend is selected explicitly before importing pyplot, because
the default backend expects a display that a server does not have.

Four reports are available: monthly, yearly, any range of months, and a detailed
one that breaks down net worth position by position and account by account.

\section{A bug worth keeping}

The first version of the doughnut chart crashed on a real month. The cause was a
category with a negative total, produced by a refund recorded as a negative
expense.

\code{matplotlib.pie} raises on slices that are zero or negative, which is
reasonable: a pie chart of a negative share has no meaning. The fix filters the
data for the chart while leaving the tables untouched:

\begin{lstlisting}[style=py]
pairs = [(l, float(v)) for l, v in zip(labels, data) if v and float(v) > 0]
if not pairs:
    ax.text(0.5, 0.5, "No data", ha="center", va="center")
    ax.axis("off")
    return _save(fig)
\end{lstlisting}

\begin{keypoint}[What the fix teaches]
The chart hides what it cannot draw; the report still shows it in the table
below. Dropping the row entirely would have been a data loss disguised as a
rendering fix. A regression test now generates a report from a month containing a
negative category.
\end{keypoint}

\section{Percentages that do not lie}

Formatting percentages sounds trivial until a real number rounds badly. A return
of 0.4\% displayed with zero decimals becomes \code{0\%}, which reads as "this
investment did nothing" when in fact it did something small.

\begin{lstlisting}[style=py]
def pct(value, decimals: int = 0, zero_decimals: int = 3) -> str:
    v = float(value or 0)
    d = zero_decimals if (v != 0 and round(v, decimals) == 0) else decimals
    return f"{v:.{d}f}%"
\end{lstlisting}

A real zero still prints as \code{0\%}. A value that only rounds to zero gets
extra precision instead. The rule is one line and it removes an entire class of
misleading output.

\section{Lazy imports, again}

The reports module pulls in the four heaviest dependencies in the project.
They are imported inside each handler, never at module level:

\begin{lstlisting}[style=py]
@router.get("/pdf/monthly")
def pdf_monthly(year: int, month: int, ...):
    from app.services.reports import generate
    return _pdf(generate.monthly_pdf(db, year, month), f"Report_{MONTHS[month]}_{year}.pdf")
\end{lstlisting}

The cost of an import inside a function is one dictionary lookup after the first
call, since Python caches modules in \code{sys.modules}. The benefit is several
seconds off every cold start, which chapter \ref{ch:operations} shows to be the
difference between a working bot and a silent one.

A test enforces the rule so it cannot rot. It runs a clean subprocess, imports
\code{app.main}, and asserts that none of the heavy modules ended up loaded.
